%%
%% Telemetry-Driven Reliability Engineering for LLM Agent Systems:
%% A Deployment Rubric from 3,780 Fault-Injection Runs
%%
%% Target: ISSRE 2026 Industry Track Full Paper (6 pages incl. references)
%% IEEE Computer Society conference format.
%% Non-anonymous submission.
%% Deadlines (AoE): Abstract 2026-06-28 (cycle 1) / 2026-07-03 (cycle 2);
%%                  Full paper 2026-07-05 (cycle 1) / 2026-07-12 (cycle 2).
%% Target: cycle 1 (abstract 2026-06-28, paper 2026-07-05).
%%

\documentclass[conference]{IEEEtran}
\IEEEoverridecommandlockouts

\usepackage{cite}
\usepackage{booktabs}
\usepackage{url}
\usepackage{enumitem}
\usepackage{amsmath}
\usepackage{multirow}
\usepackage{graphicx}
\usepackage{microtype}
\usepackage{balance}

\PassOptionsToPackage{hyphens}{url}
\setlength{\emergencystretch}{3em}
\tolerance=2000
\hfuzz=2pt
\hbadness=10000

\makeatletter
\g@addto@macro\UrlBreaks{\do\_\do\.\do\-\do\,\do\?\do\/}
\makeatother

\begin{document}

\title{Telemetry-Driven Reliability Engineering\\
for LLM Agent Systems:\\
A Deployment Rubric from 3{,}780 Fault-Injection Runs}

\author{\IEEEauthorblockN{Krishna Chaitanya Balusu}
\IEEEauthorblockA{Independent Researcher\\
San Francisco, CA, USA\\
krishnabkc15@gmail.com}}

\maketitle

\begin{abstract}
LLM agents are entering production faster than the observability
stacks that monitor them. A prior benchmark~\cite{agenttelemetry_aiware}
showed nine agent-specific span kinds detect fourteen
agent-orchestration fault classes; it stopped at a research artifact.
This paper translates the same 3{,}780-row corpus (14~faults
$\times$~6~telemetry conditions $\times$~7~frameworks $\times$~6~LLMs)
into a reliability-engineering rubric: a vendor conformance grade
card scoring seven agent SDKs A--F on out-of-the-box coverage
(detection 0.500--1.000, threshold-invariant); a blast-radius
taxonomy of the 14 faults for triage; an alert-threshold playbook
with a 7.1\% empirical FPR and a fatigue-budget calculation; a
four-week rollout pattern and runbook templates. Honest findings:
every off-the-shelf SDK ships at Grade C/D ($\sim$43\% blind spot),
the only fault firing on every production LLM
(\texttt{guardrail\_bypass}) is the one no off-the-shelf SDK detects,
and the benchmark did not measure MTTR. Artifacts are open-source.
\end{abstract}

\begin{IEEEkeywords}
LLM agents, reliability engineering, observability, AIOps.
\end{IEEEkeywords}

%% =====================================================================
\section{Introduction}
\label{sec:introduction}

Consider a composite deployment scenario synthesized from
publicly documented agent-incident patterns~\cite{mast,agenttelemetry_aiware}:
a platform team deploys a CrewAI-based pricing agent into production.
The agent enters a reasoning loop, burns a non-trivial share of the
per-task API budget before the cost-anomaly alarm fires, and the
postmortem reveals that the team's OpenTelemetry stack had no concept
of a \emph{reasoning step}---every iteration appeared as an
undifferentiated \texttt{INTERNAL} span. This pattern is not
hypothetical: it is the modal failure mode that controlled
fault-injection benchmarks on agent systems consistently
document~\cite{agenttelemetry_aiware,mast}, with the underlying
behavior (LLM agents getting stuck in repeated reasoning cycles with
cost amplification) recorded across the
1{,}600+-trace MAST corpus~\cite{mast}.

Site reliability engineering (SRE) teams adopting LLM agents inherit
a four-way deficit. The OpenTelemetry GenAI semantic
conventions~\cite{otel_genai} cover LLM invocation, tool execution,
and retrieval, but leave planning, reasoning, safety monitoring,
delegation, and memory operations as opaque internal spans;
production framework SDKs (LangChain~\cite{langchain},
CrewAI~\cite{crewai}, AutoGen~\cite{autogen},
LlamaIndex~\cite{llamaindex}, OpenAI Agents SDK, Anthropic SDK)
vary widely in how much agent-orchestration structure they expose;
microservice-era fault taxonomies (e.g.~\cite{cusick_orr,nicsdg}) do
not cover coordination and cognitive failures such as reasoning
loops, planning failure, or guardrail bypass; and the research
community has produced benchmarks and toolkits but not the
deployment-time rubrics an on-call rotation can adopt next week:
which SDK to pick, what alert threshold to set, what runbook to
write, how to score blast radius.

This paper closes the last gap. We do not propose a new benchmark, a
new span taxonomy, or a new toolkit---those were the contribution of
the AIware~2026 paper~\cite{agenttelemetry_aiware}, which we cite
throughout. Instead, we use the same 3{,}780-row benchmark as a
\emph{measurement instrument} and derive a reliability-engineering
rubric for production deployment: a vendor conformance grade card
(\S\ref{sec:conformance}), a blast-radius taxonomy of the 14 fault
classes (\S\ref{sec:blastradius}), an alert-threshold playbook
including an alert-fatigue budget and an SLO/error-budget translation
rule (\S\ref{sec:alerts}--\S\ref{sec:slo}), a four-week deployment
integration pattern with runbook templates (\S\ref{sec:deployment}),
and lessons learned (\S\ref{sec:lessons}). The artifact-by-artifact
boundary with AIware is in Table~\ref{tab:boundary}.

The audience is reliability engineers, AIOps practitioners, and
platform teams adopting LLM agents into existing
OpenTelemetry-based stacks, and the agent-SDK vendors deciding what
instrumentation to ship.

%% =====================================================================
\section{Background and Relation to Prior Work}
\label{sec:background}

\textbf{The nine agent-specific span kinds.}
The AgentTelemetry benchmark~\cite{agenttelemetry_aiware} defines nine
agent-specific span kinds derived through grounded coding of seven
production frameworks: \texttt{AGENT}, \texttt{LLM\_CALL},
\texttt{TOOL\_CALL}, \texttt{PLANNING}, \texttt{REASONING},
\texttt{RETRIEVAL}, \texttt{GUARD\_RAIL}, \texttt{DELEGATION},
\texttt{MEMORY}. Three overlap with OTel GenAI operations, one
extends the GenAI \texttt{agent} operation, and five are novel
(planning, reasoning, safety, delegation, memory).

\textbf{The 14 fault classes.}
The same work defines 14 fault classes spanning coordination
(circular delegation, agent misroute), cognitive (reasoning loop,
planning failure, hallucination), policy (guardrail bypass,
stale retrieval, memory corruption), and operational
(cost explosion, context overflow, infinite loop, tool failure,
timeout, wrong tool) categories. All 14 are documented in real
GitHub-reported incidents across six frameworks.

\textbf{The 6 telemetry conditions.}
The benchmark evaluates fault detection under six conditions of
increasing capability: no telemetry, vanilla OTel, OTel+GenAI,
OpenInference, AgentTelemetry at metadata-only privacy, and
AgentTelemetry at full-capture privacy.

\textbf{Boundary with prior work.}
The benchmark data, span taxonomy, fault taxonomy, telemetry
conditions, and toolkit were contributed
in~\cite{agenttelemetry_aiware}. This paper is a separate
contribution: we use the benchmark as a measurement instrument and
derive deployment-time engineering artifacts absent from the prior
work. Table~\ref{tab:boundary} enumerates the artifact-by-artifact
split. No tables from the prior work are imported verbatim; every
operational artifact in this paper is new.

\begin{table}[t]
\centering
\scriptsize
\caption{Artifact-by-artifact boundary with AIware 2026
(\cite{agenttelemetry_aiware}). Every operational artifact in this
paper is new; every research artifact is cited, not duplicated. A
reviewer who is suspicious of overlap can read AIware's table of
contents end-to-end and confirm none of the ``new here'' rows appears.}
\label{tab:boundary}
\begin{tabular}{@{}lccl@{}}
\toprule
\textbf{Artifact} & \textbf{AIware} & \textbf{Here} & \textbf{Why a separate treatment} \\
\midrule
9-span taxonomy           & yes & cited & Measurement primitive \\
14-fault taxonomy         & yes & cited & Measurement primitive \\
Aggregate FDR (0.612)     & yes & cited & Benchmark result \\
Ablation matrix           & yes & cited & Benchmark result \\
SWE-bench case            & yes & cited & Research validation \\
Per-framework FDR / grade & no  & Tab.~\ref{tab:conformance} & Vendor-selection rubric \\
Blast-radius S/M/L/XL     & no  & Tab.~\ref{tab:blastradius} & Triage policy \\
Alert-fatigue budget      & no  & Tab.~\ref{tab:alertbudget} & Eng-hours cost translation \\
SLO/error-budget rule     & no  & \S\ref{sec:slo}        & Reliability-currency rule \\
Four-week rollout         & no  & \S\ref{sec:deployment} & Deployment process \\
Runbook templates         & no  & \S\ref{sec:deployment} & On-call artifact \\
Postmortem fields         & no  & Tab.~\ref{tab:postmortem} & Incident-template addendum \\
\bottomrule
\end{tabular}
\end{table}

\textbf{Reliability-engineering vocabulary and benchmark scope.}
We adopt standard SRE vocabulary~\cite{sre_book, sre_workbook}
(\emph{MTTR}, \emph{MTBF}, \emph{SLO}, \emph{error budget},
\emph{blast radius}, \emph{alert fatigue}) and add
\emph{conformance} (the degree to which an implementation emits
the structural signals a downstream consumer expects). We use
controlled-benchmark FDR and FPR throughout as \emph{structural
ceilings}---what each adapter can detect when faults are present,
not how often faults will be present in production. The AIware
appendix shows organic rates in production-tier LLMs are dominated
by \texttt{missing\_guardrail} and \texttt{cost\_explosion}, with
most classes rare without injection (Lesson~3, \S\ref{sec:lessons}).

%% =====================================================================
\section{Vendor Conformance Grades for Agent SDKs}
\label{sec:conformance}

\subsection{Motivation}

A platform team picking an agent SDK today is implicitly committing
to whatever agent-orchestration structure that SDK exposes through
its native hooks. This commitment is invisible during procurement
and only surfaces at the postmortem table. We make it visible.

\subsection{Method}

For each of the seven framework adapters in
AgentTelemetry~\cite{agenttelemetry_aiware} we ran the metadata-only
condition against all 14 fault classes and all 6 mocked LLM seeds
(84 runs per framework). We compute the mean fault-detection rate
(FDR) and the count of distinct span kinds the adapter emits
natively, then assign a letter grade:
\textbf{A}~(FDR\,$\geq$\,0.95 \emph{and} 9/9 kinds);
\textbf{B}~($\geq$\,0.70 \emph{and} 7--8/9);
\textbf{C}~(0.55--0.70 \emph{or} 5--6/9);
\textbf{D}~(0.40--0.55 \emph{or} 3--4/9);
\textbf{F}~($<$\,0.40).
This weighs both \emph{detection coverage} and \emph{instrumentation
depth} so a vendor cannot mask a structural gap with
detector-specific heuristics.

\textbf{Threshold robustness.}
The central finding---one Grade A and six Grade C/D adapters---is
invariant to threshold choice within the range a reasonable rubric
author would pick. Sweeping the Grade A FDR floor over
$\{0.85, 0.90, 0.95\}$ and the kind-count floor over $\{7, 8, 9\}$
leaves the partition unchanged: \texttt{custom} is the unique A and
all six off-the-shelf adapters land at C or D in every combination.
The conformance-gap finding is not an artifact of threshold choice;
it reflects a real instrumentation-surface deficit.

\textbf{What the grade reflects.}
The grade reflects the coverage achieved by the
AgentTelemetry-shipped adapter using each framework's public hooks
(callback APIs, monkey-patch points, span handlers). It is a joint
measurement of (a)~the framework's intrinsic instrumentation
surface and (b)~the adapter's exploitation of that surface;
frameworks have not been given the opportunity to ship their own
native AgentTelemetry-compliant adapters. The rubric is therefore
a snapshot of the \emph{out-of-the-box posture} a team experiences
today, not a verdict on intrinsic framework capability; the
``Remediation path'' column estimates the kinds required to lift
each adapter one grade.

\subsection{Results}

Table~\ref{tab:conformance} reports the grade card. Only the
hand-rolled \emph{custom} adapter reaches Grade A; every off-the-shelf
SDK ships at Grade C or D today. The gap is uniform: third-party
SDKs all detect the six fault classes observable through generic
attributes (cost, tokens, tool status) and miss most of the eight
coordination, cognitive, and policy faults.

\begin{table}[t]
\centering
\scriptsize
\caption{Vendor conformance grade card. FDR is mean across 84 runs
(14 faults $\times$ 6 mocked LLMs) at metadata-only privacy. ``Span
kinds'' counts kinds emitted natively. LOC is adapter size
from~\cite{agenttelemetry_aiware}. The Grade D row falls below the
C-tier FDR floor (0.55) despite tying autogen/crewai on native kind
count, because their four kinds overlap a smaller slice of the
detection matrix.}
\label{tab:conformance}
\begin{tabular}{@{}lrrrcl@{}}
\toprule
\textbf{Framework} & \textbf{Adapter} & \textbf{Span} & \textbf{FDR} & \textbf{Grade} & \textbf{Remediation} \\
 & \textbf{LOC} & \textbf{kinds} &  &  & \textbf{path} \\
\midrule
custom        & 315 & 9/9 & 1.000 & \textbf{A} & --- (reference) \\
anthropic\_sdk & 294 & 3/9 & 0.571 & C & +5 kinds \\
autogen       & 256 & 4/9 & 0.571 & C & +4 kinds \\
crewai        & 264 & 4/9 & 0.571 & C & +4 kinds \\
openai\_sdk   & 245 & 3/9 & 0.571 & C & +5 kinds \\
langchain     & 359 & 4/9 & 0.500 & D & +4 kinds \\
llamaindex    & 280 & 4/9 & 0.500 & D & +4 kinds \\
\bottomrule
\end{tabular}
\end{table}

\subsection{Operational reading}

A team standardizing on a Grade D off-the-shelf SDK today operates
at 0.500 fault-detection coverage on the 14-class taxonomy. That is
not an algorithmic limit: every missing fault class has a defined
detection signal that requires the SDK to emit the corresponding
span kind. The remediation path is concretely sized---adding 4--5
span kinds per adapter---and the custom adapter (315 LOC, Grade A)
bounds the upper-end engineering cost. Vendor-side instrumentation
work in this range is the highest-leverage reliability investment
available to agent-SDK maintainers today.
For an SLO-bound workload, a platform team should (a)~deploy the
reference adapter for the highest-blast workloads, (b)~ship an
SDK-side patch raising the chosen vendor to Grade B+ for general
workloads, or (c)~accept the gap and tier alert policy accordingly
(\S\ref{sec:alerts}). Each option has a concrete cost in engineering
effort, vendor patching, or alert risk---this paper's contribution
is making each cost surfaceable.

%% =====================================================================
\section{Blast-Radius Taxonomy of the 14 Fault Classes}
\label{sec:blastradius}

\subsection{Method}

Incident triage requires ranking faults by operational impact, not by
detector convenience. We score each of the 14 fault classes on three
dimensions:

\begin{itemize}[nosep,leftmargin=*]
\item \textbf{Blast radius} (S/M/L/XL): scope of damage if the fault
  goes undetected over a typical detection window---direct dollar
  burn (XL), user-visible regression or safety violation (L),
  internal state drift (M), bounded transient (S). The bands are
  assigned by worst-case observed impact rather than expected impact,
  so a team that weights customer-visibility more than dollar burn
  may upgrade \texttt{memory\_corruption} and \texttt{stale\_retrieval}
  from M to L. The S/M/L/XL boundary between adjacent bands is the
  only band-pair where a reasonable SRE could disagree with our
  scoring on any of the 14 faults.
\item \textbf{Detection reliability} (Grade A/B/C/F\textsubscript{D}):
  A = detected under all conformant conditions including vanilla
  OTel; B = requires metadata-level AgentTelemetry; C = requires
  AgentTelemetry on Grade A or B framework only;
  F\textsubscript{D} = undetectable in any condition tested.
\item \textbf{Time-to-detect class} (Fast/Medium/Slow): Fast = per
  span at emit time; Medium = per trace, requires trace-level
  aggregation; Slow = cross-trace correlation or external grounding
  required. We intentionally use qualitative classes here: the
  benchmark's released TSV does not populate
  \texttt{time\_to\_root\_cause\_ms} (it is $0$ in 3{,}779 of 3{,}780
  rows; the lone $0.1$\,ms entry is an instrumentation artifact),
  so we do not report MTTR numbers. This is a known gap discussed in
  \S\ref{sec:threats}.
\end{itemize}

\subsection{Results and triage policy}

Table~\ref{tab:blastradius} presents the taxonomy. The recommended
response column gives a default triage policy a team can adapt to
their SLO; ``page'' implies immediate human attention via pager,
``ticket'' implies asynchronous queue processing, ``digest'' implies
daily review.

\begin{table}[t]
\centering
\scriptsize
\caption{Blast-radius taxonomy for the 14 fault classes. Detection
reliability anchors to the conformance grade card
(Table~\ref{tab:conformance}); A = detectable on vanilla OTel,
B = requires AgentTelemetry, C = requires Grade A/B framework.
Time-to-detect class is qualitative (TSV does not include measured
timing, see \S\ref{sec:threats}). Triage default is a
starting policy adaptable to per-tenant SLOs.}
\label{tab:blastradius}
\begin{tabular}{@{}lcccl@{}}
\toprule
\textbf{Fault class} & \textbf{Blast} & \textbf{Det.} & \textbf{TTD} & \textbf{Triage} \\
\midrule
cost\_explosion       & XL & A & Fast & page \\
circular\_delegation  & XL & B & Fast & page \\
guardrail\_bypass     & XL & C & Fast & page \\
hallucination         & L  & B & Slow & ticket \\
reasoning\_loop       & L  & C & Med  & ticket / page on recur \\
infinite\_loop        & L  & A & Fast & page \\
context\_overflow     & L  & A & Fast & page on recur \\
memory\_corruption    & L  & C & Slow & ticket \\
tool\_failure         & M  & A & Fast & ticket \\
timeout               & M  & A & Fast & ticket \\
wrong\_tool           & M  & A & Med  & ticket \\
stale\_retrieval      & M  & C & Slow & digest \\
planning\_failure     & M  & C & Med  & ticket \\
agent\_misroute       & M  & C & Med  & ticket \\
\bottomrule
\end{tabular}
\end{table}

\subsection{Operational reading}

Three observations drive the deployment plan
(\S\ref{sec:deployment}): (1)~the XL-blast faults split between
Grade A detection (\texttt{cost\_explosion}---catchable on any
stack) and Grade C detection (\texttt{guardrail\_bypass}---requires
a Grade A/B framework), making the conformance gap directly
safety-relevant; (2)~L-blast cognitive faults (hallucination,
reasoning loop, memory corruption) can compound cost before
detection fires on Grade D stacks, so \texttt{cost\_explosion} is a
useful early-warning compensating control; (3)~no fault class is
undetectable in principle---every fault has at least one
span-based signal, anchoring the conformance-investment argument
from \S\ref{sec:conformance}.

%% =====================================================================
\section{Alert Thresholds and Alert-Fatigue Budget}
\label{sec:alerts}

\subsection{False-positive baseline}

The benchmark includes 42~no-fault control runs per condition.
Under both AgentTelemetry conditions (metadata-only and
full-capture) detectors fired on 3/42 runs (FPR\,=\,0.071); under
no\_telemetry, vanilla OTel, OTel+GenAI, and OpenInference,
FPR\,=\,0.000 (those conditions detect almost nothing, including
true positives). 7.1\% FPR is the operating baseline for
metadata-level AgentTelemetry. Privacy level (metadata vs full
capture) has no effect on FDR or FPR; teams should pick metadata
for privacy compliance and accept the same alert volume.

\subsection{Threshold defaults and their sensitivity}

The benchmark's sensitivity sweep~\cite{agenttelemetry_aiware}
documents three tuning axes: \texttt{max\_retries}
(infinite\_retry, reasoning\_loop) is sensitive across 2--6
(firings 12$\to$0; production default 4 to absorb stochasticity);
\texttt{cost\_threshold} (cost\_explosion) is \emph{insensitive}
across \$0.05--\$0.15 (set by per-tenant budget policy); and
\texttt{token\_growth\_factor} (context\_overflow) is sensitive
across 1.15--2.00$\times$ (firings 21$\to$0; per-agent-type
calibration---conversational tolerates higher growth than
ReAct-style coding).

\subsection{Alert-fatigue budget}

The 7.1\% per-run FPR translates directly to operational load. For a
hypothetical platform running $N$ agent invocations per day, the
expected daily false-alert count at the defaults is
$0.071 \times N$. Table~\ref{tab:alertbudget} shows the load at three
deployment scales.

\begin{table}[t]
\centering
\scriptsize
\caption{Alert-fatigue budget at 7.1\% per-run FPR (metadata-level
AgentTelemetry, defaults). Pager budget assumes one engineer-hour per
investigated page (industry rule-of-thumb~\cite{sre_book}). At the
50K/day scale, defaults-only operation would consume 3{,}550
investigative hours/month---enough to anchor an entirely false
oncall rotation. Tiered policy (page only on Grade A detections,
ticket on Grade B/C, digest on Slow TTD) is essential at scale.
Linear-extrapolation caveat: this is a planning-grade estimate that
assumes detector firings are independent across runs and that
production traffic mirrors the benchmark's fault-class distribution.
Both assumptions are imperfect (see \S\ref{sec:threats}); teams
should treat the table as an order-of-magnitude sizing tool, not a
forecast.}
\label{tab:alertbudget}
\begin{tabular}{@{}lrrr@{}}
\toprule
\textbf{Scale} & \textbf{Runs/day} & \textbf{Exp.\ FPs/day} & \textbf{Eng-hrs/mo} \\
\midrule
small  & 1{,}000   & 71    & 71 \\
mid    & 10{,}000  & 710   & 710 \\
large  & 50{,}000  & 3{,}550 & 3{,}550 \\
\bottomrule
\end{tabular}
\end{table}

\subsection{Recommended policy}

We recommend a three-tier alert policy keyed to
Table~\ref{tab:blastradius}:

\begin{enumerate}[nosep,leftmargin=*]
\item \textbf{Page} only on XL-blast \emph{and} Grade A detection
  with the stricter threshold (e.g., cost\_explosion at 2$\times$
  budget). This bounds pager volume at roughly the cost-anomaly
  rate---typically $<$1\% of runs in healthy production.
\item \textbf{Ticket} on L-blast or Grade B/C detection. Investigated
  on the next business day with full trace evidence.
\item \textbf{Digest} on Slow-TTD or M-blast faults. Reviewed weekly
  by a designated reliability owner. Regulated workloads
  (PCI/HIPAA/SOX-bound and safety-critical agent paths) escalate one
  tier per fault class regardless of TTD, so a Slow-TTD safety class
  becomes a ticket and a Medium-TTD safety class becomes a page.
\end{enumerate}

Each tier degrades gracefully under the conformance gap from
\S\ref{sec:conformance}: a Grade D stack simply demotes Grade C
detections (rows below the line in Table~\ref{tab:blastradius}) to
``not detected'' and the corresponding fault class becomes a known
blind spot the postmortem template explicitly records (see
\S\ref{sec:deployment}).

\subsection{From conformance grade to error-budget burn: a worked example}
\label{sec:slo}

A worked example translates the conformance gap into reliability
currency. Consider a 99.5\% success-rate SLO over a 30-day window
on $N=100{,}000$ agent invocations/month: error budget is
500~failed invocations. A Grade A reference adapter detects all 14
fault classes; a Grade D off-the-shelf adapter (FDR\,$\approx$\,0.500)
detects seven (cost, tool, loop, context, timeout,
\texttt{circular\_delegation}) and misses the seven coordination,
cognitive, and policy classes that require agent-specific span
kinds---most importantly \texttt{guardrail\_bypass} (XL-blast,
FDR\,=\,0 on every off-the-shelf SDK we tested). The AIware
real-LLM appendix~\cite{agenttelemetry_aiware} reports concrete
organic detection rates across 13~production-tier models:
\texttt{missing\_guardrail} fires on 13/13, \texttt{cost\_explosion}
on 7/13, \texttt{wrong\_tool} on 2/13, \texttt{infinite\_loop} on
3/13, and the remaining ten classes register zero organic instances
in the 159-run corpus. Two consequences. \emph{First}, a Grade D
stack catches the loop/cost/tool faults that fire organically, so
the deceptive month-by-month read is ``the budget is fine.''
\emph{Second}, the one fault class that fires on every
production-tier model (\texttt{guardrail\_bypass}) is precisely the
one no off-the-shelf adapter detects: on the Grade D stack, every
guardrail bypass passes silently into the user-visible error budget,
and unlike a cost overrun (auto-cappable) the failure mode is a
safety violation the budget cannot meaningfully absorb. Closing the
\texttt{GUARD\_RAIL} span-kind gap alone---one of the five novel
kinds---moves a Grade D adapter from ``catches the noise, misses
the safety class'' to ``covers the only fault that always fires.''
The same logic applies to the ten zero-organic classes once an
upstream model regression or a chaos drill~\cite{netflix_chaos}
pushes their rate above zero.
The translation rule (conformance grade $\to$ undetected class list
$\to$ per-class organic rate $\to$ budget exposure weighted by
blast radius) is substitutable: teams plug in their own SLO,
traffic, and organic rate. This is the unit of account in which
closing the conformance gap from \S\ref{sec:conformance} becomes
funded engineering work.

%% =====================================================================
\section{Deployment Integration Pattern}
\label{sec:deployment}

\subsection{Four-week rollout}

We describe a four-week rollout pattern for adopting AgentTelemetry
inside an existing OpenTelemetry-based reliability stack (Grafana
Tempo~\cite{grafana_tempo}, Datadog, Honeycomb, or equivalent). This
pattern is designed to fail safe: each step is independently
revertible and each adds at most one new alert source. The
four-week cadence is indicative rather than prescriptive---it is
the smallest cadence that preserves the structural ordering required
by three converging bodies of deployment practice: SRE Workbook
implementing-SLOs discipline~\cite{sre_workbook} (one signal at a
time, observe before promoting to alerting), Google
progressive-rollout / canary practice~\cite{sre_book} (independent
revertibility per step, quantitative gate per promotion), and
Netflix chaos engineering~\cite{netflix_chaos} (exercise the
detection stack against known failures before relying on it). The
structural ordering---instrument first, observe baseline, enable
detectors one at a time, wire runbooks last---is what generalizes;
the wall-clock figure is illustrative and teams compress or expand
to their change-management velocity.

\noindent\textbf{Week 0 --- Inventory.}
Enumerate every agent workload in production. Score the hosting
framework on Table~\ref{tab:conformance}. Pick rollout order:
highest blast-radius workload first, lowest-grade framework last.

\noindent\textbf{Week 1 --- Bridge.}
Install the AgentTelemetry adapter for the chosen framework and
route spans through the existing OTLP collector. Per-span overhead
is on the order of $11.7\,\mu$s
median~\cite{agenttelemetry_aiware}---sub-0.006\% of typical LLM API
latencies---but every team should verify in their own stack. Gate
promotion on no regression in p99 task latency, no spike in
collector queue depth, and correct \texttt{agenttelemetry.span.kind}
tagging.

\noindent\textbf{Week 2--4 --- Detectors.}
Turn on detectors one at a time, starting with the highest-leverage
pair: cost\_explosion (XL, Grade A, Fast) and infinite\_loop (L,
Grade A, Fast). Tune thresholds against \S\ref{sec:alerts}. Observe
FPR in production traffic for 7 days before activating the next
detector; gate on FPR within 1.5$\times$ the benchmark baseline.

\noindent\textbf{Week 4+ --- Runbooks and postmortems.}
Wire each enabled detector to a runbook (template below); add an
agent-fault-class section to the postmortem template
(Table~\ref{tab:postmortem}).

\subsection{Runbook template}

Condensed runbook templates for all 14 fault classes ship in the
open-source release; two worked examples here.

\noindent\textbf{Runbook 1: reasoning\_loop.}
\emph{Trigger:} $\geq 4$ consecutive REASONING$\to$LLM\_CALL
cycles in one trace with near-identical prompt fragments.
\emph{First action:} pin to last known-good agent version if
recently changed; otherwise inspect REASONING attributes for
repeated patterns and inject a strategy-change hint or kill the
trace. \emph{Cost cap:} kill regardless of cause if trace exceeds
2$\times$ per-task budget while still looping. \emph{Postmortem
fields:} loop pattern, tool involved, prompt fragment, total cost
burn, framework grade, threshold setting.

\noindent\textbf{Runbook 2: cost\_explosion.}
\emph{Trigger:} cumulative trace cost exceeds per-task threshold
(default \$0.10). \emph{First action:} determine single-call
(unbounded prompt template---revert or apply prompt-length guard)
vs iteration-driven (cross-reference reasoning\_loop; follow
Runbook 1 if also firing). \emph{Compensating control on Grade D
stacks:} this is the primary early-warning signal for coordination
and cognitive faults that Grade C detectors miss.
\emph{Postmortem fields:} call shape, per-call token counts, agent
identity, framework grade, tenant budget tier.

\subsection{Postmortem rubric additions}

Table~\ref{tab:postmortem} lists the fields an agent-incident
postmortem should add beyond a standard SRE postmortem template
(timeline, customer impact, root cause, action items) of the kind
documented in the Google SRE Workbook~\cite{sre_workbook} and the
PagerDuty open incident-response documentation~\cite{pagerduty_ir}.

\begin{table}[t]
\centering
\scriptsize
\caption{Agent-incident postmortem fields. Capturing these per
incident creates an internal corpus that lets a team measure
their conformance-gap exposure over time and prioritize SDK-side
patches against blast-radius-weighted incident rate.}
\label{tab:postmortem}
\begin{tabular}{@{}lp{4.8cm}@{}}
\toprule
\textbf{Field} & \textbf{Purpose} \\
\midrule
Fault class & One of the 14 classes (\S\ref{sec:blastradius}) \\
Detection latency class & Fast / Medium / Slow / Not detected \\
Framework grade at time & A--F per Table~\ref{tab:conformance} \\
Blast-radius realized & Actual S/M/L/XL impact observed \\
Detector at default? & Whether alert threshold was tuned \\
Conformance gap exposed & Span kinds the SDK didn't emit \\
Remediation owner & Vendor patch, in-house adapter, accepted risk \\
\bottomrule
\end{tabular}
\end{table}

%% =====================================================================
\section{Lessons Learned and Negative Findings}
\label{sec:lessons}

The ISSRE Industry Track explicitly invites negative findings. We
surface five.

\textbf{Lesson 1: the conformance gap is industry-wide.} Every
off-the-shelf agent SDK we tested ships at Grade C or D today
(Table~\ref{tab:conformance}). This is an instrumentation deficit,
not an algorithm problem; the engineering cost to close it is
concretely bounded (4--5 additional kinds per adapter against a
245--359 LOC baseline). Until that work lands, SLO-bound workloads
on off-the-shelf SDKs operate with a documented $\sim$43\%
agent-orchestration blind spot.

\textbf{Lesson 2: metadata is enough for fault detection but
introduces a real alert-fatigue cost.} FDR and FPR are identical at
metadata-only and full-capture (both 0.612 / 7.1\% aggregate). Pick
metadata for privacy compliance; do not pay full-capture's privacy
cost expecting better detection.

\textbf{Lesson 3: controlled fault injection over-counts what
production traffic surfaces organically.} The mock benchmark reaches
FDR\,=\,1.000 with a conformance-complete adapter because faults are
injected at rate 1.0; the AIware real-LLM appendix shows organic
fault rates dominated by missing\_guardrail and cost\_explosion,
with most fault classes rare in production-tier models. Passive
monitoring alone under-validates the detection stack;
fault-injection drills (chaos engineering for
agents~\cite{netflix_chaos}) remain necessary to exercise detectors
against ground truth on a recurring cadence.

\textbf{Lesson 4: span-kind coverage has direct, fault-class-by-class
payoff.} The ablation in~\cite{agenttelemetry_aiware} shows each
span kind is the sole detection signal for at least one fault class.
Adding REASONING to a Grade D adapter immediately unlocks
reasoning\_loop detection. On the detection-coverage axis
specifically, span-kind investment has direct payoff---a property
few other reliability investments share, which typically yield
diffuse improvements rather than per-fault-class deltas.

\textbf{Lesson 5: the benchmark didn't measure time-to-detect.}
The released TSV's \texttt{time\_to\_root\_cause\_ms} column is
$0$ in 3{,}779 of 3{,}780 rows (a single row records $0.1$\,ms, an
instrumentation artifact rather than a meaningful measurement)---the
harness did not capture detection latency. This is a real
gap for operational deployment because the MTTR contribution of a
detector depends on detection latency, not just FDR. We surface this
openly rather than fabricate timing estimates; the next benchmark
iteration must instrument it.

%% =====================================================================
\section{Threats to Validity}
\label{sec:threats}

\textbf{Construct.} The conformance rubric is derived, not a
vendor-sanctioned standard, and reflects coverage on the
\cite{agenttelemetry_aiware} 14-fault taxonomy. Blast-radius
S/M/L/XL bands are qualitative; teams may re-rank under different
cost or compliance thresholds.

\textbf{Internal.} The benchmark uses deterministic mock LLM clients;
organic rates differ (Lesson~3). The 7.1\% FPR is from 42 controlled
no-fault runs per condition; production FPR may diverge.

\textbf{External.} Seven frameworks and six LLM families tested.
Conformance grades will shift as vendors update SDKs; the rubric is
versioned against the framework releases captured in the AIware
artifact. No live deployment is included; the integration pattern
(\S\ref{sec:deployment}) is \emph{designed} informed by SRE
practice~\cite{sre_book}, not \emph{measured} field outcome.

\textbf{Conclusion.} Operational claims (alert-fatigue budget,
runbook applicability, four-week rollout) are propositional---they
describe what an SRE team would do if they adopt the rubric, not
measured field outcomes.

%% =====================================================================
\section{Related Work}
\label{sec:related}

\textbf{Agent observability and benchmarks.}
AgentTelemetry~\cite{agenttelemetry_aiware} is the direct prior
work. The closest published agent-failure taxonomy is
MAST~\cite{mast}, which identifies 14 failure modes from 1{,}600+
multi-agent traces and provides the inter-annotator-validated
behavioral baseline against which our blast-radius scoring is
cross-checkable. These taxonomies characterize what fails but do
not provide deployment-time rubrics.

\textbf{LLM observability platforms.}
OpenLLMetry~\cite{openllmetry}, Langfuse~\cite{langfuse}, and
LangSmith~\cite{langsmith} are widely deployed but inherit the
OTel GenAI ceiling on agent-orchestration semantics; the
agent-specific runbook layer added here is absent from their
default configurations.

\textbf{Cloud-system reliability tooling and SRE foundations.}
ISSRE Industry Track precedent includes
Early Bird~\cite{earlybird} (failure prediction at scale), the
Global Operational Readiness Review process~\cite{cusick_orr}, and
NICSDG~\cite{nicsdg} (non-intrusive service-dependency graphs);
the contribution delta here is extending the deployment-rubric
pattern to the agent-system domain. Vocabulary and rubric
structure inherit from the SRE literature~\cite{sre_book,sre_workbook}
and the AIOps research agenda~\cite{aiops_roadmap}.

%% =====================================================================
\section{Conclusion}
\label{sec:conclusion}

Adopting LLM agents in production requires reliability-engineering
artifacts beyond benchmark evidence alone. This paper translates the
AgentTelemetry fault-injection corpus into deployment-time artifacts
that a platform team can adopt next week: a vendor conformance grade
card surfacing a $\sim$43\% blind spot across off-the-shelf SDKs,
a blast-radius taxonomy for incident triage, an alert-threshold
playbook with empirical FPR bounds and an alert-fatigue budget, a
four-week rollout pattern, and runbook templates for all 14 fault
classes. The honest negative findings---the industry-wide
conformance gap, the alert-fatigue cost of metadata-only capture,
the gap between controlled-benchmark and organic field-fault rates,
the absence of measured time-to-detect---bound what the rubric does
and does not promise. Vendors and platform teams have a concrete next
step: close the span-kind gap. The rubric, conformance card, runbook
templates, and reproduction artifacts are open-source.

%% =====================================================================
\section*{Data Availability}

The 3{,}780-row benchmark, seven framework adapters, and toolkit are
archived at \url{https://doi.org/10.5281/zenodo.20129005}
(concept DOI)~\cite{agenttelemetry_artifact} and on PyPI as
\texttt{agenttelemetry}. The conformance-grade script, blast-radius
scoring sheet, runbook templates, and postmortem rubric introduced
here are released alongside the same artifact.

\balance
\bibliographystyle{IEEEtran}
\bibliography{refs}

\end{document}
